\documentclass{article}
\usepackage[T1]{fontenc}
\usepackage{microtype}
\usepackage{graphicx}
\usepackage{booktabs}
\usepackage{hyperref}
\usepackage[preprint]{icml2026}
\usepackage{amsmath}
\usepackage{amssymb}
\usepackage{enumitem}
\usepackage{url}
\urlstyle{same}
\newcommand{\recordstatus}[1]{\noindent\textbf{Status:} #1\par}
\newcommand{\evidence}[1]{\noindent\textbf{Evidence:} #1\par}
\input{generated/metadata}
\icmltitlerunning{Dynamic Speculators}

\begin{document}
\twocolumn[
  \icmltitle{Dynamic Speculators}
  \begin{icmlauthorlist}
    \icmlauthor{Project Working Record}{project}
  \end{icmlauthorlist}
  \icmlaffiliation{project}{Internal experimental record, version \reportversion}
  \icmlcorrespondingauthor{Project record}{\href{https://github.com/AhmedZeer/dynamic-speculators-trainer/issues}{repository issues}}
  \icmlkeywords{speculative decoding, EAGLE3, LoRA, conditional generators, experimental record}
  \vskip 0.3in
]
\printAffiliationsAndNotice{}
\hypersetup{
  pdftitle={Dynamic Speculators},
  pdfsubject={Working experimental record; evidence through \evidencedate},
  pdfauthor={Project Working Record}
}

\begin{abstract}
This document records ongoing experiments with a conditional generator that
produces LoRA factors for an EAGLE3 drafter from sets of target hidden-state
summaries. It details the set encoder, shared decoder, and full-factor output
heads, and records fixed-update data scaling and the decision to retire
LoRA-bank warmup. It describes the training and serving interfaces, a supplied
four-model throughput comparison on in-distribution and out-of-distribution
workloads, activation-cache measurements, and an unresolved output-repeatability
investigation. A fixed-batch shared-adapter serving implementation is recorded
separately from its pending GPU measurements. Every experiment states its
available evidence, execution status, and missing information. The document
is extended on request and does not present final conclusions.
\end{abstract}

\section{Research Objectives and Questions}
\label{sec:questions}

We investigate a generator that maps a set of prompt summaries to one shared
adapter for a frozen EAGLE3 drafter. Separate serving studies compare the
original drafter, trained static LoRA,
shared-adapter generators and per-request generators. Checkpoint labels and
protocol changes are retained within each study rather than pooled. Historical training
ablations are recorded separately: fixed-update data scaling and the decision
to discontinue bank-based generator warmup.

The working questions are:
\begin{enumerate}[leftmargin=*]
  \item How sensitive is generator validation behavior to training-pool size
  under a fixed optimizer-update budget?
  \item How do generator checkpoints behave relative to the original drafter and
  a static trained adapter under a specified serving workload?
  \item How do recorded measurements differ between Nemotron validation prompts
  and selected Red Hat benchmark categories?
  \item What changes when one generated adapter is shared by a fixed batch of
  requests at different concurrencies?
  \item Does conditioning on more examples improve a shared adapter when
  training data and optimizer updates are matched?
  \item Do per-request adapters preserve acceptance while permitting native
  continuous admission, and does better acceptance yield higher throughput?
  \item How does LoRA rank relate to adapter variation, category separation,
  acceptance and serving token rate?
  \item What time and memory costs arise from obtaining and reusing target
  activations during training?
\end{enumerate}

These questions remain open. A numerical table records an observation under its
measurement conditions; it does not establish general superiority, statistical
significance, lossless acceleration, or completion of a planned experiment.

\subsection{Record Conventions}
Experiment identifiers remain stable as evidence is appended. We use
\emph{planned}, \emph{implemented}, \emph{partially measured}, \emph{measured}, and
\emph{investigation deferred} as execution statuses. Evidence quality is stated
separately: a measured value may still have incomplete provenance or no repetition.
Dates distinguish receipt of evidence from execution when the latter is unknown.
This version includes evidence available through \evidencedate.

\section{Method and System Description}
\label{sec:method}

\subsection{Set-Conditioned Adapter Generation}
The generator receives one summary per conditioning prompt, together with its
prompt length. Current serving uses the hidden states at the last prompt token
from the configured target layers. Prompt conditioning excludes response tokens.
For a conditioning set $S=\{(h_i,\ell_i)\}_{i=1}^{n}$, a shared encoder produces
representations $z_i$. The condition is constructed from their mean, population
standard deviation, and the logarithm of the set size:
\begin{equation}
 c=F\!\left([\operatorname{mean}_i z_i,\,
 \operatorname{std}_i z_i,\,\log(1+n)]\right).
\end{equation}
Module and layer descriptors are combined with $c$ before decoding full LoRA
factors. For each adapted projection $p$, the resulting weight is
\begin{equation}
 W'_p=W_p+\frac{\alpha}{r}B_pA_p.
\end{equation}
The current implementation adapts the drafter's V and output projections.
The set produces one adapter shared by its examples, rather than independent
adapters for each example. Set conditioning therefore differs from the number
of independently scheduled inference requests.

\subsection{Generator Architecture}
The generator is a full-factor hypernetwork: its outputs are the adapter's
actual matrices, rather than indices into a stored adapter bank. For a prompt
summary of dimension $d_h$, the shared example encoder takes
$[h_i,\log(1+\ell_i)]$ through a $(d_h+1)\!\rightarrow\!512$ linear layer,
GELU, a $512\!\rightarrow\!128$ layer, and LayerNorm. Pooling the 128-dimensional
mean and population standard deviation with $\log(1+n)$ gives 257 features.
A $257\!\rightarrow\!128\!\rightarrow\!64$ MLP with GELU produces the condition.
Pooling makes the result invariant to the ordering of conditioning examples;
the count feature distinguishes different set sizes.

Each projection has a learned 32-dimensional module descriptor, alongside a
32-dimensional layer descriptor. In the current single-layer drafter the layer
embedding has one entry. Concatenating both with the condition gives 128
features. A shared linear layer and GELU expand these to width 512, followed by
two residual MLP blocks. Each block applies LayerNorm, linear, GELU, and linear
before adding its input back.

Separate linear output heads for each projection produce
$A_p\in\mathbb{R}^{r\times d_{\mathrm{in},p}}$ and
$B_p\in\mathbb{R}^{d_{\mathrm{out},p}\times r}$.
The rank-32 generator used in the recorded logs has approximately 293 million
parameters; the large full-factor heads dominate its size. The bottleneck
condition is shared, while the factors depend on the projection descriptor.
Initially the head weights are zero, the A biases receive Kaiming initialization,
and the B biases are zero, so the initial adapter update is zero.

This structure gives one representative LoRA for a conditioning set. An ordinary
LoRA baseline instead learns one persistent pair of factors per projection.
The generator's larger trainable capacity is a confound when comparing these
approaches; matching rank alone does not match trainable parameter count.
Architecture dimensions above follow the inspected implementation defaults;
per-run configuration artifacts are needed to confirm historical overrides.

\subsection{Training Interface}
The target and base drafter remain frozen while generator parameters are
optimized through generated factors. The current recipe uses
\texttt{lk\_hybrid}, rank 32, and alpha 64. Conditioning-set sizes are sampled
from 1--8; four episodes are accumulated per optimizer update in the documented
recipe. Token-budget microbatches share the episode's factors. Context-4
validation is a separate evaluation setting. Gradients from the drafter loss
flow through the generated A/B factors into the encoder, decoder, and output
heads. Thus the generator learns adapters that serve the drafting objective
directly, without requiring a previously trained adapter as a target.

Training replays saved prompt/response token trajectories through an external
vLLM activation extractor. Only activations needed by the upcoming work are
constructed, and CPU RAM retains reusable tensors within a bounded cache.
The training pipeline does not require downloading the full activation dump.
These describe the implementation and recipe, not verified configurations of
every historical measurement.

\subsection{Serial and Fixed-Batch Serving}
Serial serving conditions one adapter on one prompt. The new fixed-batch mode
collects a closed group, obtains all prompt summaries from full target prefill,
generates one adapter, and merges it into the shared drafter before draft
prefill. The weights remain fixed until all group members finish or are aborted.
New requests cannot replace members that finish early. A subsequent group
receives its own adapter after the previous group ends.

The configured concurrency is a capacity, not a guarantee of actual batch size.
A collection timeout permits partial groups, including the final workload tail.
Original-drafter and static-LoRA baselines use the same closed admission policy
in the new sweep. Static LoRA uses a permanent merge. No GPU throughput or
acceptance result for this batching policy is included in this version.

\section{Shared Experimental Setup}
\label{sec:setup}

\subsection{Models and Data}
The documented model family is Qwen3-8B with the Red Hat EAGLE3 drafter. The
current suite compares the original drafter, a trained ordinary LoRA, and
generator snapshots labeled 2,000 and 4,000 steps. Generator data size is
reported as 24K examples. Exact training progress must ultimately come from
checkpoint metadata; labels are retained without treating them as verification.

Nemotron math, code, STEM, and chat are designated in-distribution evaluation
categories. Red Hat HumanEval, math reasoning, and summarization are designated
out-of-distribution categories in this comparison. These labels describe the
chosen evaluation sources; they do not by themselves establish absence of all
semantic overlap with training. Validation rows are excluded from optimization
by the documented trajectory membership contract.

The full-data workflow schedules six training subsets per domain, totaling
24K examples, and supports ten planned epochs. It saves Hub snapshots every
1,000 optimizer updates and at completion, with context-4 validation every 500
updates in the current recipe. The schedule is not evidence that ten epochs
have completed. This report includes only the supplied checkpoint observations.

\subsection{Metrics and Measurement Boundaries}
We keep the following quantities distinct:
\begin{itemize}[leftmargin=*]
  \item \textbf{Training and validation losses:} optimization/evaluation values
  under a specified conditioning and microbatch policy.
  \item \textbf{Offline agreement:} drafter/target argmax agreement in offline
  evaluation, not measured speculative acceptance.
  \item \textbf{Serving acceptance:} native speculative-decoding counters during
  actual generation; absent from the supplied current suite summary.
  \item \textbf{Throughput:} aggregate generated output tokens divided by
  measured wall seconds, including prefill and serving overhead.
  \item \textbf{Latency and systems timings:} request latency, first-token time,
  and stage wall time, when their artifacts are available.
\end{itemize}

\subsection{Recipe Versus Verified Run Metadata}
The independent serving implementation fixes BF16, eager execution, the V1
runner, and a pinned vLLM 0.31.0 integration. Suite defaults specify a seed,
sample limit, output budget, and warmup count. These defaults must not be
substituted for missing metadata from the supplied throughput summary.
For that summary, hardware, concurrency, sample counts, seed, workload hashes,
and exact checkpoint identities remain unverified. The cache pilot separately
reports an A100 80 GB. Neither source establishes repeated-run uncertainty.

\section{Experiments and Recorded Measurements}
\label{sec:experiments}

\subsection{GEN-003: Current Four-Model Throughput Record}
\recordstatus{Measured; supplied summary, incomplete provenance.}
\evidence{Suite \texttt{comparison-v2}; received 2026-10-22. Execution date
not supplied. Tables are generated from the frozen JSON transcription.}

\textbf{Question.} What throughput values were recorded for the four model
variants on the selected in-distribution and out-of-distribution categories?

\textbf{Setup.} The source is the suite owner's supplied
\texttt{summary.md}. Tables~\ref{tab:nemotron} and~\ref{tab:redhat} preserve all
28 reported measurements without adding older model rows. Step counts follow
the supplied labels. Correctness was not checked. This record does not include
the native per-request outputs, latency statistics, acceptance counters, or
server provenance for this suite.

\input{generated/throughput_tables}

\textbf{Record notes.} The values are recorded separately by category and
checkpoint. We do not pool categories into an overall score or infer a final
model ranking. The checkpoint pair provides observations at two reported
training stages, not a repeated or controlled causal estimate of additional
training. The two tables do not establish output quality or a target-only
decoding speedup.

\textbf{Missing evidence.} Native GuideLLM results, effective configurations,
hardware metadata, prompt fingerprints, checkpoint hashes, and repeated
measurements remain to be attached to this suite record.

\subsection{SYS-001: Cold and Warm Activation Reuse}
\recordstatus{Measured; console timings transcribed in the cache record.}
\evidence{Pilot logs received 2026-10-16; reported A100 80 GB, 400 validation
payloads evaluated in context-4 groups.}

\textbf{Question.} What stage wall times were observed before and after
activation payloads were available in the shared RAM cache?

\input{generated/cache_table}

\textbf{Record notes.} The payloads occupied 16.81 GiB. Subsequent jobs in the
same process reported initial validation passes of 19.99 and 16.86 seconds.
The first pass includes extraction and startup effects, whereas the warm pass
still performs generator/drafter evaluation. The difference is an end-to-end
observation, not an isolated cache microbenchmark or inference speedup. In a
direct-training pilot the cache reached 63.97 GiB, with 365 evictions and
69.67 GiB final process RSS. Later prefetch and summary-cache changes have no
separate GPU timing comparison in the evidence included here.

\textbf{Missing evidence.} Repeated cold/warm trials with separated extraction,
construction, transfer, and compute timings remain available follow-up work.

\subsection{GEN-002: Output Repeatability Investigation}
\recordstatus{Investigation deferred; exact sequence comparisons recorded.}
\evidence{Local correctness reports summarized in the Markdown registry,
observation dated 2026-10-21. Raw reports are not archived with this report.}

\textbf{Question.} Were target-only and speculative outputs repeatable under
the supplied prompt/settings fingerprints?

On 50 Nemotron math prompts, original versus repeated target-only decoding
matched exactly for 23/50 sequences; original versus repeated EAGLE3 decoding
matched for 50/50. Repeated target-only versus repeated EAGLE3 matched for
17/50. These are whole-sequence match counts, not token-error rates.

\textbf{Record notes.} The cause remains unresolved. Runtime environment,
numerical behavior, and decoding implementation differences have not been
isolated. Current timing experiments explicitly record correctness as unchecked
or unverified. No correctness claim is inferred from the throughput tables.

\textbf{Missing evidence.} Preserve and compare complete runtime provenance and
repeat a controlled reference experiment before making an exact-output claim.

\subsection{BATCH-001: Shared-Adapter Concurrency Sweep}
\recordstatus{Implemented; GPU measurements pending.}
\evidence{Implementation revision \texttt{c0ee0e4}; CPU lifecycle and HTTP
tests, plus an audit of the pinned vLLM hook upgrade.}

\textbf{Question.} What measurements arise when a generator conditions one
shared adapter on a fixed group at concurrency 1, 2, 4, or 8?

\textbf{Planned measurement configuration.} Reuse the frozen category prompts
for all model/concurrency variants. Record generated output tokens/s,
speculative acceptance where available, request latency, actual mean batch size,
and generator/merge/restoration costs. Keep the adapter fixed while a group
drains. Original-drafter and static-LoRA variants use matching closed groups.

\textbf{Available checks.} Tests cover packed prompt boundaries, one generator
invocation per group, restoration after the last member, delayed HTTP cleanup,
closed admission, partial groups, and suite resume across concurrency variants.
The initial broader CPU/HTTP run passed 94 tests; a later focused run after
controller refinements passed 50. These runs overlap and are not additive.

\textbf{Pending evidence.} No A100 execution, batched throughput table,
acceptance comparison, or memory measurement is supplied. CPU checks establish
neither GPU compatibility under load nor a performance benefit.

\section{Implementation and Investigation Notes}
\label{sec:implementation}

\subsection{Activation Supply and Cache Lifetime}
Training consumes local saved trajectories, with activations constructed on
demand by an external vLLM extraction server. A bounded CPU cache retains full
tensors, while prompt summaries can be reused separately. Cache capacity is
not a cap on process RSS: temporary construction buffers, model state, and
other allocations contribute additional memory. Persistent summaries,
prefetching, and concurrency controls are implementation features; only the
pilot timings in SYS-001 are treated as measurements here.

\subsection{Checkpoint Publication and Recovery}
The full-data training workflow supports local recovery and Hub snapshots
containing generator weights, optimizer state, progress, RNG state, resolved
configuration, and available provenance. Inference bundles can be downloaded
without optimizer checkpoints. Recovery capability is recorded as workflow
support; it does not demonstrate completion of the ten-epoch training plan.

\subsection{Independent Serving and Benchmarking}
The suite separates model serving from the GuideLLM client. It launches and
cleans up owned server/client process groups, exports category workloads once,
records attempt state, and can skip completed measurements with matching
artifact checksums. Model type selects an original drafter, static-LoRA bundle,
or generator bundle. Throughput can be measured without an additional
correctness-generation pass; the absence of that pass is recorded explicitly.

\subsection{Fixed-Batch Integration Boundaries}
The new scheduler leaves vacant slots empty until a group finishes. It permits
a partial group after a 50 ms collection window in the example configuration.
The runtime rejects KV preemption or incomplete initial group admission rather
than reuse draft cache entries under a different generated adapter. Larger
concurrency also increases the prefill/profile token budget and may change
GPU memory requirements. These constraints define BATCH-001
and the shared-adapter BATCH-002 protocol. GPU summary measurements are now
recorded, while detailed scheduling traces remain unavailable.

Collection, queueing, uneven request completion, and partial tails are part of
the measured serving protocol. Closed groups are not equivalent to unrestricted
continuous batching. Batch membership may vary with arrival/completion order,
even when the same prompt list is reused across model variants.

\subsection{Per-Request Multi-LoRA and Native Baseline}
The per-request controller generates factors once from complete target prefill,
retains a request-specific CPU copy, and activates only scheduled request IDs in
the draft-only native LoRA manager. Token mappings follow actual query offsets
on every draft forward, including speculative rejection and later one-token
steps. Request completion releases that request's adapter; paused request factors
remain available. The native scheduler can refill free slots without changing
other requests' adapters. CPU tests cover lifecycle and mapping invariants;
aggregate BATCH-003 counters are not a full isolation or correctness test.

Baseline launches can create an isolated package view containing the original,
checksum-verified copies of the three integration-patched vLLM files. The global
installation is left intact, and worker status can report imported source hashes.
This supports original-source EAGLE3 with native admission. Such source hashes
were not included in the supplied BATCH-003 upload, so this capability must not
be presented as independently verified provenance for those measurements.

\subsection{Bounded Profiling: Implemented, Measurements Pending}
Revision \texttt{dc872c5} adds a two-phase profiling command. An unprofiled pass
provides ordinary timing; fresh servers then collect bounded native CPU/CUDA
traces after warmup on the same frozen models and workload. The serial example
compares original EAGLE3, singleton-conditioned merged adapters and singleton
native multi-LoRA; the continuous example compares baseline and per-request
adapters at C4. Example captures use math/chat, 16 prompts per category,
128-token caps and at most 30 worker iterations.

Custom ranges cover generation, packing/loading, mapping and draft/projection
execution. Analysis records CPU inclusive/self time, GPU kernels and transfers,
union activity and gaps, and captured request occupancy. Projection ranges
include base GEMM work; isolated LoRA kernels require their own kernel rows.
The native execute-model scope includes target work and preparation rather than
providing an isolated verification timer. Unattributed activity remains unknown.
A bounded trace may miss refills or completions, and whole-workload acceptance
counters cannot supply captured-window accepted-token denominators.

CPU lifecycle tests, synthetic CUDA-event analysis, a real CPU profiler export
and mocked HTTP/GuideLLM orchestration checks were performed locally. CUDA/vLLM
profiling has not been executed locally and no user GPU traces have been supplied.
Profiler timings are diagnostic and excluded from ordinary throughput tables.
Thus recurring adapter kernels/routing remain hypotheses, not diagnosed causes.

\section{Open Questions and Planned Evidence}
\label{sec:open}

This section is a work queue rather than a conclusion. Entries are retained or
revised when new evidence arrives.

\begin{enumerate}[leftmargin=*]
  \item \textbf{Suite provenance.} Attach the native artifacts for
  \texttt{comparison-v2}; resolve exact checkpoint progress, hardware,
  concurrency, prompt identities, and effective serving settings.
  \item \textbf{Repetition.} Repeat matched serving measurements to record
  uncertainty rather than interpret small single-run differences.
  \item \textbf{Batching.} Run BATCH-001 on the intended GPU. Check actual batch
  sizes, memory, request completion/restoration, and acceptance before treating
  it as a completed performance experiment.
  \item \textbf{Reference behavior.} Continue GEN-002 when requested, separating
  target-only repeatability from adapter-related behavior.
  \item \textbf{Scaling interpretation.} Distinguish the fixed-update result
  from matched-epoch training and broader validation. GEN-001 has no repeated
  seeds or uncertainty estimates. Bank warmup is retired; GEN-004 documents
  the decision rather than scheduling further bank experiments.
  \item \textbf{Training progression.} Extend verified checkpoint observations
  as snapshots become available. Do not infer later training outcomes from
  the planned epoch schedule.
  \item \textbf{Systems costs.} Separate activation extraction, preprocessing,
  transfer, and training compute when a further timing study is commissioned.
\end{enumerate}

Future report versions should add measurements under the relevant stable
identifier, or create a new identifier if the comparison protocol changes.
Missing measurements remain missing; implementation status is not substituted
for an experimental result.

\clearpage
\appendix
\small
\section{Evidence Registry and Reproducibility}

\subsection{Source Registry}
Repository paths below identify the working records. The report's build manifest
records hashes of the source files available when this PDF was built.
\begin{description}[leftmargin=0pt]
  \item[GEN-003] \path{docs/experiments/results/GEN-003-generator-suite-throughput.md};
  the copied numeric source is \path{reporting/data/comparison-v2.json}.
  \item[SYS-001] \path{docs/experiments/eagle3_activation_ram_cache.md}; the copied
  timing source is \path{reporting/data/cache-pilot.json}.
  \item[GEN-002] \path{docs/experiments/results/GEN-002-http-token-equality.md}.
  \item[BATCH-001] \path{docs/experiments/eagle3_benchmark_suite.md} and
  \path{docs/experiments/eagle3_generator_http.md}; implementation revision
  \path{c0ee0e4b7c59ff8ff6ee8c201e591dceb0589635}.
  \item[Training] \path{docs/experiments/eagle3_generator.md},
  \path{docs/experiments/eagle3_generator_campaign_online.md}, and
  \path{docs/experiments/eagle3_generator_full_data.md}.
  \item[Datasets] \path{docs/experiments/speculator_benchmarks_inventory.md}.
\end{description}

The historical data-size ablation is preserved separately as GEN-001 in the
Markdown registry. Its qualitative screenshots and missing provenance have not
been converted into additional numerical claims or added model rows here.
Retired approaches remain outside this report's current model comparison.

\subsection{Reference Revisions and Artifact Requirements}
The documented target reference is Qwen/Qwen3-8B at
\path{b968826d9c46dd6066d109eabc6255188de91218}; the original drafter reference
is RedHatAI/Qwen3-8B-speculator.eagle3 at
\path{08610ffa01dd9f16731fe8f627b85905b6aa51c4}. The inspected Red Hat dataset
revision is \path{2ae86affa2cb97a972b7fc681dd51c04fbff083e}.
These reference revisions do not replace missing per-run checkpoint hashes.

Publishing an experiment requires its available training/evaluation/serving
commands, resolved configurations, model hashes, package/runtime details, and
source patches. Native GuideLLM outputs and workload manifests should accompany
performance records. This report contains a summary-level transcription for
GEN-003; it does not claim to include those missing run artifacts.

\noindent\begin{minipage}{\columnwidth}
\subsection{Commands and Manual Update Policy}
The shared-batch experiment is launched after updating the pinned vLLM hook:
\begingroup\footnotesize
\begin{verbatim}
python \
  scripts/apply_vllm_generator_patch.py \
  --apply
BENCH_CONFIG=examples/eagle3-lora/\
benchmark_suite.yaml
speculators benchmark suite \
  --config "$BENCH_CONFIG" \
  --output-root output/shared-batch-v1
\end{verbatim}
\endgroup
\end{minipage}

These are workflow commands, not a claim that this report's throughput tables
were produced by that new batch configuration. Actual launch commands must be
read from the associated run provenance. Use a persistent output directory on
Colab. Every server launch must pass \verb|--provenance-dir|; the suite supplies
this automatically.

Report updates occur only on explicit request. Markdown remains the working
experiment record. Before extending a report, freeze new input measurements,
update statuses and evidence dates, preserve earlier snapshots, and regenerate
tables from their numeric sources. New compiler build timestamps do not change
the evidence cutoff.

\section{Report Change History}
\textbf{Version 01, 2026-10-22.} Initial experimental record: current four-model
in-/out-of-distribution tables, cache pilot timings, deferred repeatability
investigation, and pending shared-batch GPU measurements. No final conclusions
are asserted.

\end{document}
